\chapter{Laziness and summary-producing recursion}
\source{Sample Q5, p. 5 (helper structure); Dynamic lecture (laziness context).}
Use this supplied helper. It is also available in the optional Haskell file.
\par\noindent\begin{minipage}{\linewidth}
\begin{lstlisting}
spreadMax :: [Int] -> Int -> (Int,[Int])
spreadMax [] seed = (seed,[])
spreadMax [x] seed = (x,[seed])
spreadMax (x:xs) seed = (max x rest, seed:out)
  where (rest,out) = spreadMax xs seed
\end{lstlisting}
\end{minipage}\par

\task{7.1. Evaluate the helper}
\begin{enumerate}[label=(\alph*)]
\item \code{spreadMax [] 20}
\answer{[Result.]}
\item \code{spreadMax [4] 20}
\answer{[Result.]}
\item \code{spreadMax [4,1,7] 20}
\answer{[Recursive intermediate pairs and final result.]}
\item \code{spreadMax [-8,-3,-6] 0}
\answer{[Result; does seed participate in the summary for nonempty input?]}
\end{enumerate}

\task{7.2. Explain the two outputs}
For an arbitrary finite nonempty input, describe the first and second
components. Which input determines the length of the output list? Which
component can be computed without forcing the numeric seed? Justify from the
equations rather than assuming that every argument is evaluated first.
\answer{[Component meanings and dependency explanation.]}

\task{7.3. Implement with two helper calls}
Implement a function replacing every element by the maximum of the original
list. Return \code{[]} on \code{[]}. Examples:
\code{[4,1,7] -> [7,7,7]} and \code{[-8,-3] -> [-3,-3]}.
Use exactly two calls to \code{spreadMax} for nonempty input and pair
projections or let bindings. No explicit recursion, \code{maximum}, \code{map},
\code{replicate}, folds, or other recursive list helpers.
\par\noindent\begin{minipage}{\linewidth}
\begin{codeanswer}
toMaximumTwoPass :: [Int] -> [Int]
toMaximumTwoPass = undefined  -- YOUR CODE
\end{codeanswer}
\end{minipage}\par

\task{7.4. One-call challenge}
Implement the same behavior with exactly one call to \code{spreadMax} for
nonempty input, using a recursive let/where binding. The other restrictions
remain. Handle the empty list explicitly.
\par\noindent\begin{minipage}{\linewidth}
\begin{codeanswer}
toMaximumOneCall :: [Int] -> [Int]
toMaximumOneCall = undefined  -- YOUR CODE
\end{codeanswer}
\end{minipage}\par
\answer{[Draw or describe the dependencies. Why can the summary be determined?]}

\task{7.5. Changed independent helper}
Use the helper below to return a list with each position replaced by the sum of
all input elements. Use one helper call, no other recursive list function, and
handle the empty list explicitly.
\par\noindent\begin{minipage}{\linewidth}
\begin{lstlisting}
spreadSum :: [Int] -> Int -> (Int,[Int])
spreadSum [] seed = (0,[])
spreadSum (x:xs) seed = (x+rest,seed:out)
  where (rest,out) = spreadSum xs seed
\end{lstlisting}
\end{minipage}\par
\par\noindent\begin{minipage}{\linewidth}
\begin{codeanswer}
toTotalOneCall :: [Int] -> [Int]
toTotalOneCall = undefined  -- YOUR CODE
\end{codeanswer}
\end{minipage}\par
\answer{[Explain why the changed base case matters. Predict output for [2,-1,4].]}

\task{7.6. Productivity is a dependency question}
Compare \code{let x = x + 1 in x} with
\code{let xs = 1 : xs in take 3 xs}. Explain which finite observation can be
produced and why. Do not conclude that every circular binding is safe.
\answer{[Dependency reasoning and result where one exists.]}
